\chapter*{How to read this book}
\addcontentsline{toc}{chapter}{How to read this book}
We begin with ordinary sentences. The claim \emph{Every spike belongs to a
recorded neuron} can be true, false, or ambiguous if the meaning of
\emph{belongs} is not specified. Mathematics begins by removing ambiguity:
we give names to objects, describe relations among them, and state
exactly when a claim follows from its assumptions.

A \emph{set} is a collection of objects. A \emph{relation} describes
which ordered pairs satisfy a condition. A \emph{function} assigns one output
to each admissible input. A signal is a function of time; a spike train can
be represented by a counting measure; a covariance matrix represents
pairwise relations among random variables; a neural trajectory is a
function whose values lie in a vector space.

\paragraph{How to use the numbered results.}
Results 1--92 form a prerequisite-aware catalogue of named theorems,
lemmas, inequalities, criteria, and identities. A later chapter adds nine
further results from the first inventory. The neutral word \emph{result}
is used deliberately: several entries are identities, axioms, or named
principles, not standalone theorems. The extended chapter proves a
selection; short entries are not a substitute for complete proofs.

\paragraph{Notation.} We write $x\in A$ for membership, $A\subseteq B$
for inclusion, $f:A\to B$ for a function, $\E[X]$ for expectation,
$\Cov(X,Y)$ for covariance, and $X^*$ for conjugate transpose.

\paragraph{An example before advanced mathematics.}
Let $U=\{u_1,u_2,u_3\}$ be the set of identified units and let
$T=[0,10)$ seconds. A spike event is an ordered pair
$(u,t)\in U\times T$. The relation $R\subseteq U\times T$
records which spikes occurred. A count function maps each unit
$u\in U$ to its number of spike times in $T$.
This is the progression from a sentence to sets, relations, functions,
numbers, and finally a statistical model.

\paragraph{Scientific caution.}
The repository's \texttt{arsenal/README.md} explicitly describes its
MATLAB methods as exploratory and not validated for scientific inference.
The mathematical validity of a formula is separate from the adequacy
of the assumptions required to infer physiology from data.
